\problema{House Robber III}{337}{src/02-arboles/337-house-robber-iii.cpp}{M}

El ladrón ha descubierto un nuevo vecindario para robar. Las casas en este
vecindario forman un árbol binario. La regla de oro (y única forma de alertar a
la policía) es que \textbf{no puedes robar dos casas que estén conectadas
directamente}. Es decir, si robas un nodo, no puedes robar a sus hijos
inmediatos.

Debes determinar la cantidad máxima de dinero que el ladrón puede llevarse sin
alertar a la policía.

\paragraph{Intuición.} Si estás parado en cualquier casa (nodo) tomando una
decisión, solo tienes dos opciones: o la robas, o no la robas.
\begin{itemize}
    \item \textbf{Opción 1: Robas la casa actual.} Si haces esto, la alarma de
    la policía limita tus siguientes movimientos: tienes \emph{estrictamente
    prohibido} robar las casas de los hijos directos. Así que la ganancia de
    esta rama sería el valor de la casa actual, más el dinero máximo que
    podrías obtener de los subárboles de los hijos \emph{asumiendo que no
    robaste esos hijos}.
    \item \textbf{Opción 2: No robas la casa actual.} Al mantenerte alejado de
    esta casa, la policía no sospecha. Ahora eres completamente \emph{libre}
    de hacer lo que quieras con las casas de los hijos: puedes decidir
    robarlas o no robarlas. Obviamente, elegirás la opción que te dé más
    dinero de cada hijo de forma independiente.
\end{itemize}

El problema de calcular estas dos opciones de forma independiente para cada nodo
con recursión simple es que terminarías calculando los mismos subárboles una y
otra vez, con una complejidad de $O(2^n)$.

\paragraph{Solución.} Para evitar el recálculo (y sin necesidad de usar un mapa
hash de memoización), podemos hacer que nuestra función recursiva (DFS)
devuelva \textbf{ambas} respuestas a la vez para cada subárbol. En C++, podemos
devolver un \texttt{std::pair<int, int>} con la estructura:
\texttt{\{max\_SIN\_robar, max\_ROBANDO\}}.

De esta forma, visitamos cada nodo del árbol exactamente una vez en post-orden
(visitamos a los hijos primero para tener sus resultados, y luego tomamos la
decisión para el padre).

\lstinputlisting{src/02-arboles/337-house-robber-iii.cpp}

\complejidad{$O(n)$}{$O(h)$}

\paragraph{Notas de entrevista (Amazon).} Muchos candidatos se lanzan
inmediatamente a escribir una recursión simple con memoización usando un
\texttt{unordered\_map<TreeNode*, int>}. Aunque funciona, la solución óptima es
devolver los dos estados en un par (o arreglo pequeño). Mencionar la evolución
de una solución ingenua ($O(2^n)$) a una memoizada y, finalmente, a esta de
estado compuesto ($O(n)$ de tiempo y sin la sobrecarga de un mapa hash) da una
excelente impresión en la entrevista de Amazon.
